\chapter{量子世界：光子与能级}\label{ch:g12:quantum-world}

红外取暖灯可以整晚照射你的皮肤而不留痕迹；
十分钟稀薄的山间阳光却能晒伤皮肤。这盏灯输出的
\omterm{def:g10:energy-conservation:energy}{能量}多得多——但\omterm{def:g10:energy-conservation:energy}{能量}并不是造成伤痕的原因。本章将说明，光
是一份一份到达的，只有单个\omterm{def:g10:energy-conservation:energy}{能量}足够高的\omterm{def:g12:quantum-world:photon}{光子}才能引发化学反应：
一个紫外\omterm{def:g12:quantum-world:photon}{光子}能打断一个化学键，而十亿个红外\omterm{def:g12:quantum-world:photon}{光子}却只能使其升温。
循着\omterm{def:g12:quantum-world:photon}{光子}的线索，一个世纪的探索由此展开：金属逸出电子，\omterm{def:g11:fundamental-interactions:ladder}{原子}发出谱线，
电子像波一样发生\omterm{prop:g12:light-as-wave:conditions}{干涉}。

\section{光以一份一份的形式到达：光子}

\begin{definition}[光子]\label{def:g12:quantum-world:photon}
\omterm{def:g12:oscillators-and-time:oscillator}{频率}为 $\nu$ 的光只能以整份的形式与物质交换\omterm{def:g10:energy-conservation:energy}{能量}，
这些整份称为\emph{光子}\index{光子}，每个光子的\omterm{def:g10:energy-conservation:energy}{能量}为
\[
E = h\nu = \frac{hc}{\lambda}, \qquad h = \qty{6.63e-34}{J.s},
\]
其中 $h$ 为\emph{普朗克常量}\index{普朗克常量}（普朗克，
1900；爱因斯坦，1905），$\lambda$ 为\omterm{def:g12:light-as-wave:lightwave}{真空中的波长}。光束的
光强表示每秒光子的\emph{数目}；每个光子的
\omterm{def:g10:energy-conservation:energy}{能量}仅由\omterm{def:g12:oscillators-and-time:oscillator}{频率}决定。
\end{definition}

\begin{notation}[电子伏特]\label{not:g12:quantum-world:ev}
\omterm{def:g11:fundamental-interactions:ladder}{原子}物理中常用\emph{电子伏特}\index{电子伏特}，
$\qty{1}{eV} = \qty{1.60e-19}{J}$：一个\omterm{def:g11:fundamental-interactions:elementary}{元电荷}
经过 \qty{1}{V} 的电势差时获得的\omterm{def:g10:energy-conservation:energy}{能量}（见\cref{ch:g11:electric-gravitational-fields}）。便捷的近似关系为：
$E \,(\unit{eV}) \approx 1240/\lambda\,(\unit{nm})$.
\end{notation}

\begin{omfigure}
\begin{tikzpicture}[font=\small, x=1.48cm]
  \draw[-Stealth] (7.55,0) -- (-0.55,0);
  \draw[-Stealth, omDef] (4.9,-0.85) -- (2.3,-0.85)
    node[midway, below] {每个光子的能量随之增大};
  \foreach \x/\l/\e/\n in {0/{\qty{1}{pm}}/{\qty{1.2}{MeV}}/{$\gamma$},
      1.2/{\qty{0.1}{nm}}/{\qty{12}{keV}}/{X},
      2.4/{\qty{100}{nm}}/{\qty{12}{eV}}/{紫外},
      3.6/{\qty{550}{nm}}/{\qty{2.3}{eV}}/{可见光},
      4.8/{\qty{10}{\micro m}}/{\qty{0.12}{eV}}/{红外},
      6.0/{\qty{1}{cm}}/{\qty{0.12}{meV}}/{微波},
      7.2/{\qty{3}{m}}/{\qty{0.4}{\micro eV}}/{无线电}}{
    \draw (\x,0.07) -- (\x,-0.07);
    \node[above, omExo] at (\x,0.42) {\n};
    \node[above, font=\scriptsize] at (\x,0.10) {\l};
    \node[below, font=\scriptsize, omDef] at (\x,-0.10) {\e};}
\end{tikzpicture}

{\small 电磁波家族的统一刻度
（见\cref{ch:g12:light-as-wave}）：\omterm{def:g12:mechanical-waves:wavelength}{波长}（上方），每个\omterm{def:g12:quantum-world:photon}{光子}的\omterm{def:g10:energy-conservation:energy}{能量}
$hc/\lambda$（蓝色）：从\omterm{def:g10:energy-conservation:energy}{能量}仅为微\omterm{def:g12:nuclear-energy:units}{电子伏特}量级的无线电\omterm{def:g12:quantum-world:photon}{光子}，到
\unit{MeV} 伽马\omterm{def:g12:quantum-world:photon}{光子}（见\cref{ch:g11:nucleus-radioactivity}）——然而一支
\qty{1.0}{mW} 的激光笔每秒发出 $\num{3e15}$ 个\omterm{def:g12:quantum-world:photon}{光子}：日常
光线的流动如倾泻的沙子般平滑。}
\end{omfigure}

\section{光电效应}

让光照射真空中洁净的金属板，电子便会逸出——这就是
\emph{光电效应}\index{光电效应}。其规律与经典物理相悖。

\begin{proposition}[实验揭示的事实]\label{prop:g12:quantum-world:facts}
对每种金属，都存在一个\emph{截止频率}\index{截止频率} $\nu_0$：
\begin{enumerate}
  \item \omterm{def:g12:oscillators-and-time:oscillator}{频率}低于 $\nu_0$ 时，无论光多强，都没有电子逸出；
  \item \omterm{def:g12:oscillators-and-time:oscillator}{频率}高于 $\nu_0$ 时，无论光多弱，电子都会立即逸出；
  \item 光强增大时，每秒逸出电子的\emph{数目}增加，
        但不改变其最大\omterm{def:g11:mechanical-energy:kinetic}{动能}；\omterm{def:g12:oscillators-and-time:oscillator}{频率}增大时，$E_{k,\max}$ 增大。
\end{enumerate}
这些现象都无法用连续波解释：按连续波的图像，任何颜色的光都应使电子逐渐积累\omterm{def:g10:energy-conservation:energy}{能量}——
哪怕是微弱的光，也要数月才能积累（见\cref{exo:g12:quantum-world:14}）——而且
光强本应决定电子的\omterm{def:g10:energy-conservation:energy}{能量}：三条预言全部落空。
\end{proposition}
\admitted

\begin{theorem}[爱因斯坦光电效应方程]\label{thm:g12:quantum-world:einstein}
\index{光电效应方程}要使电子逸出金属表面，至少需要金属的\emph{逸出
功}\index{逸出功} $W_0$。
一个\omterm{def:g12:quantum-world:photon}{光子}被一个电子整体吸收，于是
\[
h\nu = W_0 + E_{k,\max}, \qquad \text{因此 } \nu_0 = \frac{W_0}{h}.
\]
\end{theorem}

\begin{proof}
\omterm{def:g10:energy-conservation:energy}{能量}核算：输入 $h\nu$，逸出时消耗 $W_0$，其余为\omterm{def:g11:mechanical-energy:kinetic}{动能}——
\omterm{def:g11:mechanical-energy:kinetic}{动能}至多为 $h\nu - W_0$，从较深处逸出的电子\omterm{def:g11:mechanical-energy:kinetic}{动能}更小。\omterm{def:g12:oscillators-and-time:oscillator}{频率}低于 $\nu_0$ 时
单个\omterm{def:g12:quantum-world:photon}{光子}的\omterm{def:g10:energy-conservation:energy}{能量}不足以支付\omterm{thm:g12:quantum-world:einstein}{逸出功}，无论\omterm{def:g12:quantum-world:photon}{光子}数目多少。
\end{proof}

\begin{definition}[遏止电压]\label{def:g12:quantum-world:stopping}
要测量 $E_{k,\max}$，可加反向\omterm{def:g11:circuits-and-power:voltage}{电压}使电子减速
（见\cref{ch:g11:electric-gravitational-fields}）；\emph{遏止
电压}\index{遏止电压} $U_s$ 是使电流恰好为零的最小\omterm{def:g11:circuits-and-power:voltage}{电压}：
\omterm{def:g11:mechanical-energy:kinetic}{动能}最大的电子恰好不能到达另一极，因此 $E_{k,\max} = e\,U_s$。
\end{definition}

\begin{omfigure}
\begin{tikzpicture}
\begin{axis}[omaxis, width=8.8cm, height=5.4cm,
    xlabel={$\nu$ ($10^{14}\,\unit{Hz}$)}, ylabel={$E_{k,\max}$ (\unit{eV})},
    xmin=0, xmax=13.4, ymin=-3.1, ymax=3.2,
    xtick={2,4,6,8,10,12}, ytick={-2,-1,1,2,3}]
  \addplot[omExo, dashed, domain=0:5.55, samples=2] {0.414*x-2.3};
  \addplot[omThm, thick, domain=5.55:12.4, samples=2] {0.414*x-2.3};
  \addplot[only marks, mark=*, mark size=1.5pt, omDef] coordinates
    {(7,0.60) (8,1.01) (9,1.43) (10,1.84) (11,2.26)};
  \node[below, font=\small] at (axis cs:5.55,-0.08) {$\nu_0$};
  \node[right, font=\small, omExo] at (axis cs:0.15,-2.55) {$-W_0$};
   \node[font=\small, omThm, rotate=22] at (axis cs:10.3,1.35) {斜率 $= h$};
\end{axis}
\end{tikzpicture}

{\small 测得的 $E_{k,\max} = eU_s$ 与\omterm{def:g12:oscillators-and-time:oscillator}{频率}（钠）的关系：
是一条斜率为 $h$ 的直线，与坐标轴交于 $\nu_0$；
外推（虚线）可得截距 $-W_0$，位于 $\nu = 0$ 处。}
\end{omfigure}

\begin{method}[测量普朗克常量]\label{met:g12:quantum-world:measure-h}
\begin{enumerate}
  \item 用经过已知\omterm{def:g12:mechanical-waves:wavelength}{波长}\omterm{def:g11:color-light-sources:subtractive}{滤光片}的光照射光电管；
        算出每个 $\nu = c/\lambda$，记录每个\omterm{def:g12:quantum-world:stopping}{遏止电压}。
  \item 作 $E_{k,\max} = eU_s$ 随 $\nu$ 变化的图像：爱因斯坦预言
        图线为一条直线 $h\nu - W_0$。
  \item 读取斜率：它\emph{正是} $h$；截距给出 $\nu_0$ 和 $-W_0$。
\end{enumerate}
\end{method}

\section{能级：原子为何发出谱线}

\begin{definition}[能级]\label{def:g12:quantum-world:levels}
\omterm{def:g11:fundamental-interactions:ladder}{原子}中束缚电子的\omterm{def:g10:energy-conservation:energy}{能量}具有
\emph{量子化}\index{量子化}的特征：只能取离散的
\emph{能级}\index{能级} $E_1 < E_2 < \dots < 0$，
\omterm{def:g10:energy-conservation:energy}{能量}以自由电子为零点计为负值。最低的一级是\emph{基态}\index{基态}，
其余各级称为\emph{激发态}\index{激发态}；
达到 $0$ 即为
\emph{电离}\index{电离}。
\end{definition}

\begin{proposition}[氢原子的能级阶梯]\label{prop:g12:quantum-world:hydrogen}
氢\omterm{def:g11:fundamental-interactions:ladder}{原子}的\omterm{def:g12:quantum-world:levels}{能级}为
\[
E_n = -\frac{\qty{13.6}{eV}}{n^2}, \qquad n = 1, 2, 3, \dots
\]
故 $E_1 = \qty{-13.6}{eV}$，$E_2 = \qty{-3.40}{eV}$，$E_3 =
\qty{-1.51}{eV}$，$E_4 = \qty{-0.85}{eV}$，\omterm{def:g12:quantum-world:levels}{能级}向 $0$ 处越来越密集。
\end{proposition}
\admitted

\begin{remark}[能级阶梯的来源]\label{rem:g12:quantum-world:ladder-origin}
这些数值的来历属于量子力学，留待大学分册讨论——
德布罗意公式亦然：电子的\omterm{def:g12:quantum-world:debroglie}{物质波}
（见下文\cref{def:g12:quantum-world:debroglie}）必须在\omterm{def:g11:fundamental-interactions:ladder}{原子核}周围自身闭合，
只有某些\omterm{def:g12:mechanical-waves:wavelength}{波长}能够吻合——如同吉他弦，一个
\omterm{def:g11:fundamental-interactions:ladder}{原子}也只有一组离散的\omterm{def:g12:sound-acoustics:pitch-loudness}{音调}。
\end{remark}

\begin{theorem}[跃迁的光子]\label{thm:g12:quantum-world:transition}
\omterm{def:g11:fundamental-interactions:ladder}{原子}只能通过发射或吸收一个\omterm{def:g12:quantum-world:photon}{光子}来改变\omterm{def:g12:quantum-world:levels}{能级}，该\omterm{def:g12:quantum-world:photon}{光子}恰好携带两\omterm{def:g12:quantum-world:levels}{能级}之差：
由\omterm{def:g12:quantum-world:levels}{能级} $p$ 跃迁到较低\omterm{def:g12:quantum-world:levels}{能级} $n$ 时，发出
$h\nu = E_p - E_n$；由 $n$ 升至 $p$ 时吸收同样的\omterm{def:g10:energy-conservation:energy}{能量}。
\end{theorem}

\begin{proof}
\omterm{def:g10:energy-conservation:energy}{能量}守恒，逐份核算：\omterm{def:g12:quantum-world:photon}{光子}是光唯一认可的
“货币”，\omterm{def:g11:fundamental-interactions:ladder}{原子}的“账目”必须平衡。
\end{proof}

\begin{remark}[线状光谱的解释]\label{rem:g12:quantum-world:lines}
这就是\cref{ch:g10:light-spectra}中所许诺的答案：热气体只发出
仅属于其自身阶梯的\omterm{def:g12:oscillators-and-time:oscillator}{频率} $(E_p - E_n)/h$——亮线构成的“条形码”；
而冷气体恰好从\omterm{def:g10:light-spectra:white}{白光}中吸走这些\omterm{def:g12:oscillators-and-time:oscillator}{频率}，
留下对应的暗吸收线。每种元素都有自己的
阶梯，也就有了各自的“条形码”——我们正是借此解读星光。
\end{remark}

\begin{example}[氢原子的可见光谱线]\label{ex:g12:quantum-world:balmer}
终态为 $n = 2$ 的跃迁落在可见光范围内。取 $E \,(\unit{eV}) =
1240/\lambda\,(\unit{nm})$：$3 \to 2$：\qty{1.89}{eV}，\qty{656}{nm}（红色）；
$4 \to 2$：\qty{2.55}{eV}，\qty{486}{nm}（蓝绿色）；$5 \to 2$：
\qty{2.86}{eV}，\qty{434}{nm}（紫色）。跃迁到 $n = 1$ 的，
\omterm{def:g10:energy-conservation:energy}{能量}均超过 \qty{10}{eV}，全部是\omterm{def:g10:light-spectra:wavelength}{紫外线}。
\end{example}

\begin{omfigure}
\begin{tikzpicture}[font=\small, y=1cm]
  \foreach \y/\n/\e in {0/1/-13.6, 4.50/2/-3.40, 5.33/3/-1.51}{
    \draw[thick] (0,\y) -- (4.2,\y);
    \node[left] at (0,\y) {$n=\n$};
    \node[right, font=\scriptsize] at (4.2,\y) {\qty{\e}{eV}};}
  \draw[thick] (0,5.63) -- (4.2,5.63) node[left,at start,yshift=-2.5pt] {$n=4$};
  \draw[thick] (0,5.76) -- (4.2,5.76) node[left,at start,yshift=3.5pt] {$n=5$};
  \draw[dashed] (0,6.0) -- (4.2,6.0) node[left, at start] {$\infty$}
    node[right, font=\scriptsize] {\qty{0}{eV}（已电离）};
  \foreach \x/\ys/\col/\l in {1.0/5.33/omThm/656, 1.9/5.63/omDef/486,
      2.8/5.76/violet/434} \draw[-Stealth, \col, thick] (\x,\ys) --
    (\x,4.50) node[midway, left, font=\scriptsize] {\qty{\l}{nm}};
  \draw[-Stealth, omExo, thick] (3.7,4.50) -- (3.7,0)
    node[midway, right, font=\scriptsize] {\qty{122}{nm}（紫外）};
\end{tikzpicture}

{\small 氢\omterm{def:g11:fundamental-interactions:ladder}{原子}\omterm{def:g12:quantum-world:levels}{能级}阶梯（按比例绘制）：\omterm{def:g10:light-spectra:wavelength}{可见光谱}线均终止于
$n = 2$；跃迁到\omterm{def:g12:quantum-world:levels}{基态}则属深紫外。}
\end{omfigure}

\begin{remark}[激光概述]\label{rem:g12:quantum-world:laser}
\omterm{def:g10:energy-conservation:energy}{能量}恰为 $E_p - E_n$ 的\omterm{def:g12:quantum-world:photon}{光子}经过一个\emph{\omterm{def:g12:quantum-world:levels}{激发态}}\omterm{def:g11:fundamental-interactions:ladder}{原子}时，可以触发它
发出一个全同\omterm{def:g12:quantum-world:photon}{光子}——\omterm{def:g10:energy-conservation:energy}{能量}相同、方向相同、步调一致。若保持
高\omterm{def:g12:quantum-world:levels}{能级}\omterm{def:g11:fundamental-interactions:ladder}{原子}多于低\omterm{def:g12:quantum-world:levels}{能级}\omterm{def:g11:fundamental-interactions:ladder}{原子}，一个\omterm{def:g12:quantum-world:photon}{光子}便会引发雪崩式的复制：
\emph{受激辐射}\index{受激辐射}——而\omterm{def:g11:color-light-sources:lamps}{激光器}不过如此。
\end{remark}

\section{物质波}

\begin{proposition}[电子的干涉]\label{prop:g12:quantum-world:interference}
让电子逐个通过双缝
（见\cref{ch:g12:light-as-wave}）：每个电子都只落成一个清晰的点——
一个粒子。但当成千上万个电子累积时，这些点便描绘出\omterm{def:g12:light-as-wave:fringes}{干涉条纹}：
每个电子单独地以波的形式穿过\emph{两条}缝。晶体
同样能使电子束发生\omterm{def:g12:light-as-wave:diffraction}{衍射}（戴维孙与革末，1927 年）。
\end{proposition}
\admitted

\begin{omfigure}
\begin{tikzpicture}[font=\small]
  \foreach \o in {0,4.0,8.0} \draw[omExo] (\o,0) rectangle +(3.4,2.2);
  \foreach \x/\y in {0.62/1.71, 1.68/0.42, 2.31/1.95, 1.05/1.22,
      2.98/0.66, 1.74/1.58, 0.38/0.35, 2.40/0.98, 1.62/1.02, 0.97/1.86,
      2.95/1.60, 1.71/0.20, 1.12/0.61, 2.33/0.31}
    \fill (\x,\y) circle (0.8pt);
  \foreach \c/\f in {0.40/610, 1.05/410, 1.70/530, 2.35/350, 3.00/470}
    \foreach \y in {0.25,0.62,...,2.00}
      \fill ({4.0+\c+0.14*sin(\f*\y)},\y) circle (0.8pt);
  \foreach \x/\y in {4.72/1.42, 5.38/0.52, 6.06/1.12, 6.66/1.72,
      5.44/1.88, 7.28/0.78} \fill (\x,\y) circle (0.8pt);
  \foreach \c/\f in {0.40/610, 1.05/410, 1.70/530, 2.35/350, 3.00/470}
    \foreach \y in {0.12,0.27,...,2.10}
      \fill ({8.0+\c+0.11*sin(\f*\y)},\y) circle (0.8pt);
  \foreach \c/\f in {0.40/260, 1.05/340, 1.70/205, 2.35/290, 3.00/385}
    \foreach \y in {0.20,0.50,...,2.05}
      \fill ({8.0+\c+0.06*sin(\f*\y+40)},\y) circle (0.8pt);
  \node[below] at (1.7,-0.05) {$\sim 30$ 个电子};
  \node[below] at (5.7,-0.05) {$\sim 200$}; \node[below] at (9.7,-0.05) {$\sim 2000$};
\end{tikzpicture}

{\small 单个电子只留下一个点；成千上万个电子则形成条纹。下一个点落在何处
是随机的——类似于单个\omterm{def:g11:fundamental-interactions:ladder}{原子核}的衰变（见\cref{ch:g11:nucleus-radioactivity}）
——但波决定了\emph{概率}：亮纹处，落点的可能性大；暗纹处，
几乎从不落点。个体的机遇由精确的波所支配，这正是量子物理的核心。}
\end{omfigure}

\begin{definition}[德布罗意波长]\label{def:g12:quantum-world:debroglie}
每个\omterm{def:g12:newtons-laws:momentum}{动量}为 $p = mv$ 的粒子都以\emph{物质波}\index{物质波}的形式传播，
其\emph{德布罗意波长}\index{德布罗意波长}
为
\[
\lambda = \frac{h}{p} = \frac{h}{mv}.
\]
\end{definition}
\admitted

\begin{example}[为何宏观物体不会衍射]\label{ex:g12:quantum-world:never}
经 \qty{100}{V} \omterm{def:g11:circuits-and-power:voltage}{电压}加速的电子，其 $E_k = \qty{1.60e-17}{J}$，
$p = \sqrt{2m_e E_k} = \qty{5.4e-24}{kg.m/s}$，故 $\lambda \approx
\qty{0.12}{nm}$——恰为晶体中\omterm{def:g11:fundamental-interactions:ladder}{原子}的间距，因而晶体充当
它的双缝。一个 \qty{58}{g} 的网球以 \qty{25}{m/s} 运动：$\lambda =
6.63\times10^{-34}/1.45 \approx \qty{4.6e-34}{m}$，比\omterm{def:g11:fundamental-interactions:ladder}{原子核}约小 $10^{19}$ 倍
——任何狭缝都不可能使它显示\omterm{def:g12:light-as-wave:diffraction}{衍射}条纹。$h$ 极小，
以至量子颗粒性只在\omterm{def:g11:fundamental-interactions:ladder}{原子}尺度上才会显现。
\end{example}

\begin{example}[电子显微镜]\label{ex:g12:quantum-world:microscope}
\omterm{def:g12:light-as-wave:diffraction}{衍射}限制了分辨率，使人无法分辨远小于所用\omterm{def:g12:mechanical-waves:wavelength}{波长}的细节
（见\cref{ch:g12:light-as-wave}）：可见光的分辨极限约为 \qty{0.5}{\micro m}。
经 \qty{10}{kV} \omterm{def:g11:circuits-and-power:voltage}{电压}加速的电子，其 $\lambda \approx \qty{12}{pm}$，
比可见光短数万倍：电子显微镜能够观察病毒、细胞
结构，甚至单个\omterm{def:g11:fundamental-interactions:ladder}{原子}。\omterm{def:g12:quantum-world:debroglie}{物质波}为物理学提供了一只新的眼睛。
\end{example}

\section{练习}

\begin{exercise}[$\star$]\label{exo:g12:quantum-world:1}
一支绿色激光笔发出\omterm{def:g11:work-of-force:power}{功率} \qty{1.0}{mW}、\omterm{def:g12:mechanical-waves:wavelength}{波长} \qty{530}{nm} 的光。求一个
\omterm{def:g12:quantum-world:photon}{光子}的\omterm{def:g10:energy-conservation:energy}{能量}（分别以\omterm{def:g11:work-of-force:work}{焦耳}和\omterm{def:g12:nuclear-energy:units}{电子伏特}表示），再求每秒的\omterm{def:g12:quantum-world:photon}{光子}数。
\end{exercise}

\begin{exercise}[$\star$]\label{exo:g12:quantum-world:2}
一座调频电台以 \qty{100}{MHz} 发射。求其\omterm{def:g12:quantum-world:photon}{光子}\omterm{def:g10:energy-conservation:energy}{能量}（以
\unit{eV} 为\omterm{def:g10:orders-of-magnitude:unit}{单位}），并与可见光\omterm{def:g12:quantum-world:photon}{光子}比较：为何无线电的颗粒性无法被探测？
\end{exercise}

\begin{exercise}[$\star$]\label{exo:g12:quantum-world:3}
铯的\omterm{thm:g12:quantum-world:einstein}{逸出功}为 $W_0 = \qty{1.9}{eV}$。求其\omterm{prop:g12:quantum-world:facts}{截止频率}与
对应的\omterm{def:g12:mechanical-waves:wavelength}{波长}。\omterm{def:g12:mechanical-waves:wavelength}{波长}为 \qty{633}{nm} 的红色激光能否使铯发射电子——
若能，则光电子的最大\omterm{def:g11:mechanical-energy:kinetic}{动能}为多少？
\end{exercise}

\begin{exercise}[$\star$]\label{exo:g12:quantum-world:4}
利用 $E_n = -13.6/n^2$ \unit{eV}，求氢\omterm{def:g11:fundamental-interactions:ladder}{原子}
$3 \to 2$ 跃迁所发出\omterm{def:g12:quantum-world:photon}{光子}的\omterm{def:g10:energy-conservation:energy}{能量}与\omterm{def:g12:mechanical-waves:wavelength}{波长}，并判断其颜色。
\end{exercise}

\begin{exercise}[$\star$]\label{exo:g12:quantum-world:5}
由本课程中 $E_{k,\max}$ 与 $\nu$ 的图像，读出钠的\omterm{prop:g12:quantum-world:facts}{截止频率}
与\omterm{thm:g12:quantum-world:einstein}{逸出功}（以 \unit{eV} 为\omterm{def:g10:orders-of-magnitude:unit}{单位}），该图线的斜率代表什么常量？
\end{exercise}

\begin{exercise}[$\star\star$]\label{exo:g12:quantum-world:6}
\omterm{def:g12:mechanical-waves:wavelength}{波长}为 \qty{400}{nm} 的光照射到钾（\omterm{thm:g12:quantum-world:einstein}{逸出功} $W_0 = \qty{2.3}{eV}$）上。
求\omterm{def:g12:quantum-world:photon}{光子}\omterm{def:g10:energy-conservation:energy}{能量}、光电子的最大\omterm{def:g11:mechanical-energy:kinetic}{动能}（分别以 \unit{eV} 和\omterm{def:g11:work-of-force:work}{焦耳}表示），
以及\omterm{def:g12:quantum-world:stopping}{遏止电压}。
\end{exercise}

\begin{exercise}[$\star\star$]\label{exo:g12:quantum-world:7}
用\omterm{def:g12:oscillators-and-time:oscillator}{频率}高于\omterm{prop:g12:quantum-world:facts}{截止频率}的光照射光电管。在\omterm{def:g12:oscillators-and-time:oscillator}{频率}固定、光强加倍时，(a) 电流、
(b) \omterm{def:g12:quantum-world:stopping}{遏止电压}各如何变化？当\omterm{def:g12:oscillators-and-time:oscillator}{频率}升高时又如何？
哪一事实与波动图像相矛盾？请说明理由。
\end{exercise}

\begin{exercise}[$\star\star$]\label{exo:g12:quantum-world:8}
氢\omterm{def:g11:fundamental-interactions:ladder}{原子}从\omterm{def:g12:quantum-world:levels}{基态}\omterm{def:g12:quantum-world:levels}{电离}所需的最小\omterm{def:g12:quantum-world:photon}{光子}\omterm{def:g10:energy-conservation:energy}{能量}是多少？若从
$n = 2$ \omterm{def:g12:quantum-world:levels}{能级}\omterm{def:g12:quantum-world:levels}{电离}呢？求两种情形的截止\omterm{def:g12:mechanical-waves:wavelength}{波长}，并指出它们所属的\omterm{def:g10:light-spectra:white}{光谱}区域。
\end{exercise}

\begin{exercise}[$\star\star$]\label{exo:g12:quantum-world:9}
让\omterm{def:g10:light-spectra:white}{白光}穿过一罐冷氢气（所有\omterm{def:g11:fundamental-interactions:ladder}{原子}均处于\omterm{def:g12:quantum-world:levels}{基态}）。
试解释：透射\omterm{def:g10:light-spectra:white}{光谱}在可见光范围内\emph{没有}暗线，
而热恒星的\omterm{def:g10:light-spectra:white}{光谱}中却出现了巴耳末吸收线。
\end{exercise}

\begin{exercise}[$\star\star$]\label{exo:g12:quantum-world:10}
求下列粒子的\omterm{def:g12:quantum-world:debroglie}{德布罗意波长}：(a) 经
\qty{100}{V} \omterm{def:g11:circuits-and-power:voltage}{电压}加速的电子；(b) 质量为 \qty{0.10}{g}、
速率为 \qty{1.0}{m/s} 的苍蝇。将两者分别与\omterm{def:g11:fundamental-interactions:ladder}{原子}大小（$\sim \qty{e-10}{m}$）作比较，并作出结论。
\end{exercise}

\begin{exercise}[$\star\star$]\label{exo:g12:quantum-world:11}
暗适应的人眼对约 $90$ 个 \qty{505}{nm} 光的\omterm{def:g12:quantum-world:photon}{光子}产生响应。
求该\omterm{def:g10:energy-conservation:energy}{能量}，并将其与一只
质量为 \qty{2.5}{mg}、速率为 \qty{0.50}{m/s} 的蚊子的\omterm{def:g11:mechanical-energy:kinetic}{动能}作比较，并解释该比值。
\end{exercise}

\begin{exercise}[$\star\star\star$]\label{exo:g12:quantum-world:12}
某光电管在 \qty{405}{nm} 光照下给出 $U_s = \qty{1.19}{V}$，而在 $U_s =
\qty{0.40}{V}$（\qty{546}{nm}）时，由这两个数据测得 $h$，再求
以 \unit{eV} 表示的\omterm{thm:g12:quantum-world:einstein}{逸出功}。本章中的哪种金属是阴极？
\end{exercise}

\begin{exercise}[$\star\star\star$]\label{exo:g12:quantum-world:13}
某氢灯在 \qty{486}{nm} 和 \qty{434}{nm} 处有谱线。将
每条谱线换算为\omterm{def:g12:quantum-world:photon}{光子}\omterm{def:g10:energy-conservation:energy}{能量}，并借助匹配
\omterm{def:g12:quantum-world:levels}{能级}之差（$E_n = -13.6/n^2$ \unit{eV}）判断两条谱线分别对应的跃迁。
\end{exercise}

\begin{exercise}[$\star\star\star$]\label{exo:g12:quantum-world:14}
经典\omterm{def:g12:mechanical-waves:delay}{延迟}估算：强度为 \qty{1.0e-6}{W/m^2} 的微弱光
照射到钾上（$W_0 = \qty{2.3}{eV}$）。若一个\omterm{def:g11:fundamental-interactions:ladder}{原子}只能收集
穿过其横截面（半径 \qty{1.0e-10}{m}）的\omterm{def:g10:energy-conservation:energy}{能量}，则
它需要多长时间才能积累到 $W_0$？并与观测到的\omterm{def:g12:mechanical-waves:delay}{延迟}（小于 \qty{e-9}{s}）作比较。
\end{exercise}

\begin{exercise}[$\star\star\star$]\label{exo:g12:quantum-world:15}
在电子显微镜中，电子经 \qty{5.0}{kV} 的\omterm{def:g11:circuits-and-power:voltage}{电压}加速。
求其\omterm{def:g12:quantum-world:debroglie}{德布罗意波长}，并与绿光（\qty{550}{nm}）的\omterm{def:g12:mechanical-waves:wavelength}{波长}作比较，
并结合\cref{ch:g12:light-as-wave}中的\omterm{def:g12:light-as-wave:diffraction}{衍射}知识，解释为何这能提高分辨率。
\end{exercise}

\section{问题：太阳能电池板物理实验}

\begin{problem}[{周末问题——太阳能电池板物理实验：屋顶每秒吸收的
$10^{21}$ 个\omterm{def:g12:quantum-world:photon}{光子}，能测定\omterm{def:g12:quantum-world:photon}{普朗克常量}的光电管，
以及逐条读取谱线的氢灯——一个周末，一个 $h$}]
\label{pb:g12:quantum-world:1}
你们班将在一个周末内对学校屋顶的太阳能电池板进行仪器测量。
数据：正午太阳辐照度为 \qty{1.0e3}{W/m^2}；太阳光的平均\omterm{def:g12:mechanical-waves:wavelength}{波长}为
\qty{550}{nm}；电池板面积 \qty{1.0}{m^2}，\omterm{def:g10:energy-conservation:efficiency}{效率} $20\%$；当日的
日照量相当于 \qty{5.0}{h} 的满日照。

\medskip
\noindent\textbf{第一部分——\omterm{def:g12:quantum-world:photon}{光子}核算。}
\begin{enumerate}
  \item 求一个平均太阳光\omterm{def:g12:quantum-world:photon}{光子}的\omterm{def:g10:energy-conservation:energy}{能量}，分别用\omterm{def:g11:work-of-force:work}{焦耳}和\omterm{def:g12:nuclear-energy:units}{电子伏特}表示。
  \item 推算正午时每秒射到电池板上的\omterm{def:g12:quantum-world:photon}{光子}数。
  \item 求正午时电池板的电\omterm{def:g11:work-of-force:power}{功率}，再求
        一天内所产生的\omterm{def:g10:energy-conservation:forms}{电能}，并以 \unit{kW.h} 为\omterm{def:g10:orders-of-magnitude:unit}{单位}。
  \item 电池板的电流看起来非常平稳。利用第 2 题，
        解释为何\omterm{def:g12:quantum-world:photon}{光子}逐个到达的过程是不可察觉的。
  \item 太阳光中\omterm{def:g10:light-spectra:wavelength}{红外线}携带的\omterm{def:g10:energy-conservation:energy}{能量}多于\omterm{def:g10:light-spectra:wavelength}{紫外线}，但
        只有\omterm{def:g10:light-spectra:wavelength}{紫外线}能晒伤皮肤（化学键能约 $\sim \qtyrange{3}{4}{eV}$）。从\omterm{def:g12:quantum-world:photon}{光子}角度解释。
\end{enumerate}

\medskip
\noindent\textbf{第二部分——用实验套件中的光电管测量 $h$。}
实验套件中的光电管阴极材料未知，用\omterm{def:g11:color-light-sources:subtractive}{滤光片}得到的数据如下：
\begin{center}\small
\begin{tabular}{lcccc}
$\lambda$ (\unit{nm}) & 365 & 405 & 436 & 546\\
$U_s$ (\unit{V}) & 1.50 & 1.16 & 0.94 & 0.37
\end{tabular}
\end{center}
\begin{enumerate}[resume]
  \item \omterm{def:g12:quantum-world:stopping}{遏止电压}测量的是什么？请写出 $U_s$、$\nu$、$W_0$、$h$ 之间的关系。
  \item 将四个\omterm{def:g12:mechanical-waves:wavelength}{波长}换算为\omterm{def:g12:oscillators-and-time:oscillator}{频率}（\omterm{def:g10:orders-of-magnitude:unit}{单位}为 $10^{14}\,\unit{Hz}$）。
  \item 为什么作 $eU_s$ 对 $\nu$ 的图像应为一条直线？
        检验各数据点是否共线。
  \item 由两端的数据点求斜率，即测得的 $h$。
  \item 求阴极的\omterm{thm:g12:quantum-world:einstein}{逸出功}（\unit{eV}）、\omterm{prop:g12:quantum-world:facts}{截止频率}
        和截止\omterm{def:g12:mechanical-waves:wavelength}{波长}。
  \item 太阳光中哪一部分能使该
        阴极射出电子——\omterm{def:g10:light-spectra:white}{光谱}中哪一半对它毫无用处？
  \item 将你测得的 $h$ 与 \qty{6.63e-34}{J.s} 相比较，相对误差为百分之几？
\end{enumerate}

\medskip
\noindent\textbf{第三部分——标定用氢灯。}
实验套件的\omterm{def:g12:mechanical-waves:wavelength}{波长}刻度用氢灯进行校核，
其\omterm{def:g12:quantum-world:levels}{能级}为 $E_n = -13.6/n^2$ \unit{eV}。
\begin{enumerate}[resume]
  \item 求 $E_2$、$E_3$、$E_4$ 和 $E_5$。
  \item 求跃迁 $3 \to 2$、$4 \to 2$、$5 \to 2$ 对应的\omterm{def:g12:quantum-world:photon}{光子}\omterm{def:g10:energy-conservation:energy}{能量}。
  \item 推算这三条谱线的\omterm{def:g12:mechanical-waves:wavelength}{波长}，并说明各谱线的颜色。
  \item 为何氢灯发出的是\emph{谱线}，而非连续的
        彩虹状\omterm{def:g10:light-spectra:white}{光谱}？用一句话指出其中的关键思想。
  \item 该氢灯中的氢能否吸收 \qty{500}{nm} 的光？请利用\omterm{def:g12:quantum-world:levels}{能级}阶梯说明理由。
\end{enumerate}

\medskip
\noindent\textbf{第四部分——检查电池片：电子波。}
电池片中的微裂纹远小于光学显微镜的
\qty{0.5}{\micro m} 分辨极限；实验室的扫描电子显微镜工作\omterm{def:g11:circuits-and-power:voltage}{电压}为 \qty{10}{kV}。
\begin{enumerate}[resume]
  \item 求电子的速率，并用经典公式检验其是否小于 $0.25\,c$。
  \item 求电子的\omterm{def:g12:newtons-laws:momentum}{动量}和\omterm{def:g12:quantum-world:debroglie}{德布罗意波长}。
  \item 分辨率与\omterm{def:g12:mechanical-waves:wavelength}{波长}相关：求相对于绿光
        （\qty{550}{nm}）的\omterm{def:g10:signals-and-waves:oscilloscope}{增益}，并结合本周末得到的两个数值——你测得的 $h$
        和这一\omterm{def:g10:signals-and-waves:oscilloscope}{增益}，作出结论。
\end{enumerate}
\end{problem}
